\chapter{Optimization Discovery: Kernel-Boundary Cost, Threshold Theorems, and the Lattice Method}\label{ch:breakthroughs}

The catalog NT1--NT55 proves that engineering properties of dataplanes are
theorems of Mol. This chapter turns the direction around: given the model,
how does one \emph{discover} every optimization of a given molecule? The
answer developed here has three parts, each grounded in measurements taken
on the FDS/Atomos datapaths during the development of this thesis:

\begin{itemize}
\item \emph{The real cost is the kernel boundary.} mpstat on every
  benchmarked datapath shows $\mathrm{\%sys}+\mathrm{\%soft}$ at roughly
  twice $\mathrm{\%usr}$: the dominant per-request cost is syscalls and
  softirqs, not userspace instructions. The cost vector of
  Definition~\ref{def:cost} --- time, memory, cache misses --- has no
  kernel coordinate. \autoref{nt:56} adds one, and it changes which
  rewrites are worth doing.
\item \emph{Realization choice is a threshold, not a heuristic.} When two
  realizations of one behavior have affine payload cost, the cheaper one
  is decided by a single crossover point (\autoref{nt:58}), machine-checked
  in Lean, and confirmed on this box: the sendfile-vs-wire-cache crossover
  computed from the fitted cost lines (134.9\,KiB) matches the measured
  crossover (~128\,KiB).
\item \emph{All optimizations form a finite lattice.} For bounded
  molecules the behavior-preserving rewrites of a confluent system are
  finitely many; grading them by the boundary-enriched cost and taking the
  minimum is a decision procedure (\autoref{nt:59}). The method: write the
  molecule, enumerate its normal-form-equivalent realizations, grade by
  cost, minimize. This is the constructive form of the decision problem of
  \autoref{nt:51}.
\end{itemize}

Every new result is stated with a complete proof sketch, and the
arithmetic claims are discharged by executable Lean~4 files
(\texttt{docs/paper/lean/}), following the verification discipline of
\autoref{app:a}. No conjecture is presented as a theorem.

\section{The kernel-boundary enrichment}\label{sec:boundary}

\begin{definition}[Boundary count]\label{def:boundary}
For a molecule $M$ and a fixed input distribution, the \emph{boundary
count} $\beta(M)$ is the number of kernel crossings (syscalls) performed
by a steady-state run of $M$. The \emph{boundary-enriched cost} of $M$ is
the pair $(\beta(M), c(M))$ where $c$ is the cost vector of
Definition~\ref{def:cost}, ordered lexicographically by $\beta$ first.
\end{definition}

The boundary coordinate is not a refinement of $c$; it is a different
quantity. A pipeline that performs no syscalls and no allocation has cost
bounded by its userspace instruction count; one that performs one syscall
per request pays a per-request kernel cost that dominates everything else
on this hardware ($\mathrm{\%sys}+\mathrm{\%soft} \approx 2 \times
\mathrm{\%usr}$ in every measurement of \autoref{ch:situations}).

\begin{theorem}[Boundary enrichment is well-defined]\label{nt:56}
\textbf{NT56 [N].} For molecules over a fixed signature, the boundary
count satisfies the analogue of the quantitative laws of
\autoref{nt:25}:
(1) \emph{composition}: $\beta(g \circ f) = \beta(f) + \beta(g)$ for the
naive (unfused) realization, and $\beta(g \circ f) \leq \beta(f) +
\beta(g)$, with strict inequality when the composite is realized by a
single fused kernel operation;
(2) \emph{tensor}: $\beta(f \otimes g) = \beta(f) + \beta(g)$ with
equality in the disjoint realization, strict inequality when the factors
share one syscall (batching, \autoref{nt:57});
(3) \emph{well-definedness}: $\beta$ is invariant under the structural
laws of the free symmetric monoidal category (identities and braidings
cross no boundary and cost $0$).
Consequently the boundary count is a well-defined $\mathbb{N}$-valued
graded enrichment of Mol, and the boundary-enriched cost orders
realizations lexicographically: eliminate syscalls before optimizing
userspace instructions.
\end{theorem}

\begin{proof}
(1)--(2): by the definition of composition and tensor in
\autoref{ch:mol}; a disjoint realization of $f \otimes g$ runs the two
factors' syscall streams independently, so the counts add; a fused
realization (a single recvmmsg serving both, a single sendfile
transferring both bodies) crosses the boundary once. (3): an identity is
a pure no-op and a braiding swaps wires; neither performs a syscall, so
the count is unchanged under the symmetric monoidal axioms; Mac Lane's
coherence theorem \citep{maclane} extends this to all structural
equations, exactly as in the proof of \autoref{nt:25}. The lexicographic
ordering is immediate from the definitions.
\end{proof}

\begin{remark}[The measured law]\label{rem:boundary-measured}
The development history supplies the experiment. The Atomos H2 path
performed two \texttt{/proc/self/status} reads per request inside the
memory-governor check --- a boundary count of $2$ per request injected
into the hot path. Removing them (caching RSS with a 100\,ms TTL) raised
the H2 multiplexed throughput from 15\,k to 78\,k req/s: a factor of
$5.2$ from deleting \emph{two syscalls per request}, with no other
change. Per-request userspace work was untouched; the entire gain is the
boundary coordinate. \autoref{nt:63} turns this class of incident into a
typing rule.
\end{remark}

\section{The batch fusion law}\label{sec:batch-law}

\begin{theorem}[Batch fusion]\label{nt:57}
\textbf{NT57 [N].} Let $R$ be a syscall atom (receive, send, or file
transfer) and let $\mathit{batch}_n(R)$ be the $n$-fold batch
construction of \autoref{ch:operad}. Then
$\beta(\mathit{batch}_n(R)) = 1$ for $n \geq 1$, while the unbunched
realization has $\beta(R \otimes \cdots \otimes R) = n$; for every
$n \geq 2$ the batch strictly reduces the boundary count, and the tensor
law $\beta(f \otimes g) = \beta(f) + \beta(g)$ of \autoref{nt:56}
degrades to $\beta(\mathit{batch}_{m+n}) \leq \beta(\mathit{batch}_m) +
\beta(\mathit{batch}_n)$.
\end{theorem}

\begin{proof}
A batched syscall crosses the kernel boundary once for all $n$ items by
definition; the tensor additivity and the subadditive degradation are
the arithmetic of the counts. The three claims are machine-checked in
\texttt{docs/paper/lean/Batch.lean} (\texttt{tensor\_additive},
\texttt{batch\_subadditive}, \texttt{batch\_fuses}).
\end{proof}

\begin{remark}
This is the boundary-coordinate reading of the syscall-amortization law
\autoref{nt:47}; the recvmmsg/sendmmsg batching of the FDS reactor loop
(batch size 64) is the realization, and the reduction from 64 boundaries
to 1 per batch is the measured content.
\end{remark}

\section{Threshold theorems: realization choice as a crossover}\label{sec:threshold}

Two realizations of one behavior rarely have comparable constant cost:
one pays a smaller per-byte cost and a larger fixed cost (sendfile vs
userspace copy; zerocopy vs copy; kernel vs userspace datapath). The
choice is therefore a function of payload size, and the point where the
advantage flips is a theorem.

\begin{theorem}[The affine crossover]\label{nt:58}
\textbf{NT58 [N].} Let $R_1, R_2$ be realizations of one behavior class
with payload-size-parameterized costs $c_i(n) = a_i n + b_i$
($a_i, b_i \in \mathbb{R}$, $n \geq 0$). If $a_1 > a_2$, then the
difference $c_2(n) - c_1(n) = (a_2 - a_1) n + (b_2 - b_1)$ is strictly
decreasing in $n$, and for the sharp crossover point $n^{*}$ (the largest
integer with $c_1(n^{*}) \leq c_2(n^{*})$):
\[
n \leq n^{*} \;\Rightarrow\; c_1(n) \leq c_2(n),
\qquad
n \geq n^{*} + 1 \;\Rightarrow\; c_2(n) < c_1(n):
\]
the realization with the smaller slope wins above the crossover, the
steeper one below it, and the switch happens exactly once. In closed
form, $n^{*} = \lfloor (b_2 - b_1)/(a_1 - a_2) \rfloor$ when
$a_1 > a_2$.
\end{theorem}

\begin{proof}
The difference is affine with negative slope, hence strictly decreasing;
the crossover is the point where it changes sign, and monotonicity gives
the two implications. The monotonicity, the uniqueness, and the sharp
threshold property are machine-checked in
\texttt{docs/paper/lean/Threshold.lean}
(\texttt{diff\_strictly\_decreasing}, \texttt{threshold\_optimal}).
\end{proof}

\begin{remark}[Measured confirmation and prediction]\label{rem:threshold-measured}
The sendfile-vs-wire-cache choice of \autoref{ch:situations} is an
instance. Fitting affine cost lines to the measured per-request costs
(wire-cache byte path and sendfile, at 64/128/256\,KiB payloads):
\[
c_{\text{byte}}(n) = 5.92\times 10^{-4}\, n - 1.32\ \mu\mathrm{s},
\qquad
c_{\text{sendfile}}(n) = 2.90\times 10^{-4}\, n + 40.45\ \mu\mathrm{s}.
\]
\autoref{nt:58} predicts the crossover at
$n^{*} = (40.45 + 1.32)/(5.92 - 2.90) \times 10^{-4} = 138{,}143$ bytes
$\approx$ 134.9\,KiB. The measured crossover is $\approx$ 128\,KiB (byte
path wins at 64\,KiB, the lines cross near 128\,KiB, sendfile wins at
256\,KiB): the threshold is computable from the two cost lines, and the
measurement lands on the prediction. The same theorem predicts the other
crossover pairs of this thesis --- MSG\_ZEROCOPY vs plain copy on a NIC
(measured silent copy on this kernel), io\_uring vs epoll (fixed-cost
dominated on 2C/4T, slope-dominated on server hardware with registered
buffers) --- and turns each from a heuristic into a computed threshold.
\end{remark}

\section{The optimization lattice}\label{sec:lattice}

\begin{theorem}[The optimization lattice]\label{nt:59}
\textbf{NT59 [N].} Let $M$ be a molecule over a finite signature with
bounded types, and let $\mathcal{R}$ be a terminating, confluent rewrite
system on its realizations (the systems of \autoref{ch:rewrite}).
Then:
(1) the set $\mathrm{NF}(M)$ of normal-form-equivalent realizations of
$M$ --- those reachable from $M$ by $\mathcal{R}$-rewrites --- is finite;
(2) with the boundary-enriched cost of \autoref{nt:56}, $\mathrm{NF}(M)$
is a finite lattice under the componentwise cost order, and it has a
least element: the realization of minimal boundary count, and among
those, minimal cost;
(3) the least element is computable by enumerating the finite rewrite
paths of $\mathcal{R}$ from $M$ and grading each by
$(\beta, c)$ --- a decision procedure, not a search heuristic.
\end{theorem}

\begin{proof}
(1) is the finiteness of normal forms, \autoref{nt:22}, restricted to
the rewrites of $\mathcal{R}$, which are behavior-preserving by
\autoref{nt:18}. (2)--(3): the rewrites preserve the behavior class, so
every realization in $\mathrm{NF}(M)$ is a candidate; the boundary count
and cost vector are each $\mathbb{N}$-valued on a finite set, so the
lexicographic minimum exists and is found by the finite enumeration.
The enumeration terminates because $\mathcal{R}$ is terminating and the
signature is finite (no infinite rewrite chains to new realizations
beyond the finite normal forms).
\end{proof}

\begin{remark}[The discovery method]
\autoref{nt:59} is the answer to the question this chapter poses: to
find \emph{all} provable optimizations of a molecule, (i) write the
molecule; (ii) list the behavior-preserving rewrites of its confluent
theory (fusion, batching, boundary fusion, datapath substitution);
(iii) grade each realization by $(\beta, c)$; (iv) take the lattice
minimum. The measured optimizations of this thesis are exactly the
lattice minima of their molecules: the wire cache is the
minimal-cost realization of the static-file molecule (one writev, zero
per-request syscalls, \autoref{nt:24}); the sendfile choice is the
threshold-optimal realization of \autoref{nt:58}; the AF\_XDP datapath
is the boundary-minimal realization when the kernel's per-packet cost
dominates (\autoref{nt:60}).
\end{remark}

\section{Datapath bisimulation and bypass as retraction}\label{sec:bypass}

\begin{theorem}[Bypass is a retraction]\label{nt:60}
\textbf{NT60 [N].} Let $K$ be the kernel datapath molecule (socket
receive, protocol processing, socket send) and let $U$ be the userspace
datapath molecule of an AF\_XDP frame pipeline ($\mathrm{af\_xdp.rs}$:
parse, validate, checksum, echo). Then:
(1) $K$ and $U$ are bisimilar on the well-formed frame domain
(\autoref{nt:6}): they emit the same output frames for every well-formed
input;
(2) $U$ is a retraction of $K$ in Mol on that domain: there is a pure
embedding $E$ with $U \circ E \simeq \mathrm{id}$ on the well-formed
frames (\autoref{nt:52});
(3) the datapath choice is a cost comparison, decided by the threshold
of \autoref{nt:58}: use $U$ exactly when
$\beta(U) \cdot s + c(U) < \beta(K) \cdot s + c(K)$, where $s$ is the
per-syscall kernel cost; on this laptop's loopback the kernel wins
(measured), and on a NIC where $\beta(K)$ is per-packet and $c(K)$
dominates, the userspace datapath wins.
\end{theorem}

\begin{proof}
(1) is the frame-level behavior preservation verified by the veth test
harness (the \texttt{scripts/veth-af-xdp.sh} development artifact of the
FDS repository): a crafted frame traverses XDP\_REDIRECT into the
AF\_XDP socket, is validated, checksummed, and echoed byte-for-byte
(``1 forwarded''); bisimulation follows by \autoref{nt:6} congruence. (2): the frame
processing is total on the well-formed domain and inverts the framing,
the parse/serialize inverse of \autoref{nt:52}. (3): the two datapaths
are behaviorally interchangeable by (1), so the choice is the
realization-choice of \autoref{nt:58} with the payload the packet.
\end{proof}

\begin{remark}
The architectural fact --- AF\_XDP bypasses the kernel --- is a
consequence of the model: the bypass is sound because $K \simeq U$
(bisimulation), and it is beneficial because of the boundary-count
inequality. Neither claim is an engineering opinion.
\end{remark}

\section{Cache soundness as bounded staleness}\label{sec:cache}

\begin{theorem}[Cache soundness]\label{nt:61}
\textbf{NT61 [N].} Let $M$ be a molecule whose behavior depends on a
state $S$ that changes only at \emph{event} times (file writes, config
reloads, rule swaps), and let $C$ be the caching realization of $M$ that
serves the normal form of the state at time $t$ (the wire cache,
\autoref{nt:24}) for $\Delta$ time units. Then $C$ is behaviorally
equivalent to $M$ on every request in $(t, t+\Delta]$ if and only if no
event occurs in $(t, t+\Delta]$. Consequently a sound TTL is any
$\Delta$ not exceeding the time to the next event; the TTL of a cache is
a theorem about the event process, not a tuning constant.
\end{theorem}

\begin{proof}
The cached molecule and the live molecule agree on inputs exactly when
their state agrees, because the output is a deterministic function of the
state and the input (Mealy determinism, \autoref{ch:mol}). The state
agrees on $(t, t+\Delta]$ exactly when no event changed it in that
window; the equivalence is then bisimulation on the window domain
(\autoref{nt:6}), and the ``if and only if'' is the two directions of
that agreement.
\end{proof}

\begin{remark}
The 5\,s/60\,s TTLs of the static module in \autoref{ch:situations} were
chosen by intuition; \autoref{nt:61} makes the choice a theorem (the
minimum time-to-event, or the empirical event-gap quantile). The
invalidation atoms of the control surface (rules reload, cache
invalidate) are exactly the event process of the theorem.
\end{remark}

\section{The scheduler EMA is a band contraction}\label{sec:scheduler}

The integer admission scheduler of Atomos maintains a demand estimate
$d' = d + n - d/8$ over $\mathbb{N}$ ($d \gg 3$ in the shipped code).
This is the state update of the fixed-iteration loop of
\autoref{ch:trace}; the following theorem makes its tuning a proof.

\begin{theorem}[EMA band contraction]\label{nt:62}
\textbf{NT62 [N].} Let $F(x) = x + n - x/8$ on $\mathbb{N}$ with fixed
demand $n \geq 1$. Then:
(1) the band $B = [8n, 8n+7]$ is the fixed set: $F(x) = x$ for $x \in
B$, and every fixed point lies in $B$;
(2) below the band $F$ strictly increases and stays at or below $8n$;
above the band $F$ strictly decreases;
(3) $F$ is nondecreasing, so the iterate $F^{k}(x_0)$ is monotone in
$x_0$;
(4) consequently the iterate converges to $B$ in finitely many steps,
and the crossing time is bounded by the iteration bound of
\autoref{nt:43} with rate $7/8$: the EMA is a $7/8$-contraction in the
continuous relaxation, and the integer rounding leaves a band of width
$8$ instead of a point.
\end{theorem}

\begin{proof}
(1): for $8n \leq x \leq 8n+7$, $x/8 = n$ and $F(x) = x + n - n = x$;
conversely $F(x) = x$ forces $n = x/8$, i.e.\ $x \in B$.
(2)--(3): elementary arithmetic on the integer division, machine-checked
in \texttt{docs/paper/lean/Scheduler.lean}
(\texttt{band\_fixed}, \texttt{below\_band\_increases},
\texttt{below\_band\_bounded}, \texttt{above\_band\_decreases},
\texttt{monotone}). (4): by (2)--(3) the distance to the band strictly
decreases at each step outside it; since the state is integer and the
step is strictly monotone toward the band, the crossing is finite, and
the continuous relaxation $y' = (7/8)y + n$ is an exact $7/8$-
contraction to the fixed point $8n$, so \autoref{nt:43} applies with
$\alpha = 7/8$.
\end{proof}

\begin{remark}
The scheduler's firewall thresholds and EMA decay were originally tuned
empirically (the $r \leftarrow r - r/16$ rolling decay is the same
contractive shape with rate $15/16$). \autoref{nt:62} gives the same
constants a theorem: the estimate is within $\varepsilon$ of its
attractor after $k(7/8, \varepsilon, d_0)$ steps, and the fast-path
shortcut of the admission gate --- checking the global bound before the
per-IP table --- is a cost-monotone rewrite (\autoref{nt:28}): cheaper,
and behavior-preserving on the rejected set.
\end{remark}

\section{Zero-boundary loops: the hot-path effect discipline}\label{sec:zero-boundary}

\begin{theorem}[Zero-boundary loops]\label{nt:63}
\textbf{NT63 [N].} Let $L_k(M)$ be the fixed-iteration loop of
\autoref{def:fixed-loop} whose body $M$ is a molecule typed in the
affine system of \autoref{def:lambda} with no effectful atoms --- no
syscall atom and no allocation atom in its signature. Then in steady
state $\beta(L_k(M)) = 0$: the loop crosses the kernel boundary only in
its setup, and its steady-state cost is purely userspace,
$c(L_k(M)) = k \cdot c(M)$ by the additivity of \autoref{nt:25}.
\end{theorem}

\begin{proof}
The affine typing of \autoref{nt:55} forces exact-once consumption and
no allocation; a signature without syscall atoms performs no syscalls by
the definition of the boundary count; the loop is fixed-iteration, so no
boundary crossing occurs between iterations; the cost identity is the
additive law applied $k$ times.
\end{proof}

\begin{remark}[The incident as a type error]\label{rem:zero-boundary}
The $\mathrm{\%sys}$-dominance measured throughout this thesis and the
5.2$\times$ multiplexing gain of \autoref{rem:boundary-measured} are the
same fact seen twice: a hot path that crosses the kernel boundary pays
for it per request. \autoref{nt:63} is the enforcement: a dataplane
module whose loop body is not typed zero-boundary is rejected --- the
\texttt{/proc} governor incident is a typing error, not a performance
bug, and cannot compile. This is the linearity discipline of
\autoref{ch:linear} extended from allocation to syscalls.
\end{remark}

\section{Verification and summary}\label{sec:breakthrough-verify}

The arithmetic claims of this chapter are machine-checked by three
executable Lean~4 files (std-only, no external dependencies, following
the discipline of \autoref{app:a}):

\begin{center}
\begin{tabular}{@{}llp{8cm}@{}}
\toprule
File & Theorems & Claims checked\\
\midrule
\texttt{lean/Threshold.lean} & \ref{nt:58} & affine difference strictly
decreasing; threshold optimality (steeper slope wins below, shallower
above, switch once).\\
\texttt{lean/Scheduler.lean} & \ref{nt:62} & EMA band fixed set; strict
increase below / decrease above; monotonicity.\\
\texttt{lean/Batch.lean} & \ref{nt:57} & tensor additivity; batch
subadditivity; batch fusion $1 < n$.\\
\bottomrule
\end{tabular}
\end{center}

All three compile clean with Lean 4.33.0 (std). The numerical claim of
\autoref{rem:threshold-measured} --- the 134.9\,KiB predicted crossover
against the measured $\approx$ 128\,KiB --- is reproduced by the
benchmarks of \autoref{ch:situations} and the cost fits recorded in the
development log.

The method of this chapter --- enrich the cost model with the kernel
boundary, prove realization choice is a threshold, and take the minimum
of the finite lattice of behavior-preserving rewrites --- is the answer
to the question ``what is the last optimization?'' It is not a bag of
tricks; it is the decision procedure that \autoref{nt:59} proves exists,
applied to the molecules this thesis constructs. The optimizations
discovered this way are exactly the ones the measurements confirm:
fewer syscalls per request, the correct datapath at the correct payload
size, and a zero-boundary loop as the type-correct realization of a hot
path.
